\begin{theorem}
	Dane jest ciało \( \mathbb{F} \) charakterystyki \( p \) i liczba \( q = p^k \) dla pewnego \( k \). Jeśli wielomian \( f(X) = X^q - X \) ma \( q \) pierwiastków, to stanowią one \( q \)-elementowe podciało \( \mathbb{F} \).
\end{theorem}
\begin{proof}
	Niech \( A = \{a \in \mathbb{F}: f(a) = 0\} = \{a \in \mathbb{F}: a^q = a\} \) będzie zbiorem pierwiastków. Wystarczy pokazać, że zbiór \( A \) jest zamknięty na działania, bo inne aksjomaty ciała (łączność, przemienność, rozdzielność, etc.) wynikają z faktu, że pracujemy w~\( \mathbb{F} \), które od początku je spełniało.
	\begin{itemize}
		\item Elementy neutralne: \( 0, 1 \in A \)
		\item Element \( -1 \in A \): \\
		      \( (-1)^q = -1 \) (również dla \( p = 2 \))
		\item Elementy odwrotne: \\
		      \( \pars{a^{-1}}^q \cdot a = \pars{a^{-1}}^q \cdot a^q = 1 \), czyli \( \pars{a^{-1}}^q = a^{-1} \in A \)
		\item Zamkniętość na mnożenie: \\
		      \( (a \cdot b)^q = a^q \cdot b^q = a \cdot b \) \\
		      W szczególności, jeśli \( a \in A \), to \( -a = -1 \cdot a \in  A \), stąd mamy elementy odwrotne w dodawaniu.
		\item Zakmniętość na dodawanie: \\
		      Najpierw udowodnimy, że w \( \mathbb{F} \) zachodzi \( (a + b)^p = a^p + b^p \) dla dowolnych \( a, \ b \). \\
		      Wiemy, że:
		      \[
			      (a + b)^p = a^p + \ldots + {p \choose i}a^{p-i}b^i + \ldots + b^p
		      \]
		      Wszystkie współczynniki \( p \choose i \) dla \( 1 \leq i < p \) są podzielne przez \( p \), bo licznik jest równy \( p! \) \linebreak a mianownik jest niepodzielny przez \( p \).
		      Zatem w ciele charakterystyki \( p \) wszystkie składniki oprócz \( a^p, \ b^p \) są równe \( 0 \).
		      Korzystając z tego, dostajemy:
		      \[
			      (a + b)^q = ((a + b)^p)^{p^{k-1}} = (a^p + b^p)^{p^{k-1}} = (a^{p^2} + b^{p^2})^{p^{k-2}} = \ldots = a^{p^k} + b^{p^k} = a^q + b^q = a + b
		      \]
	\end{itemize}
\end{proof}

\subsection{Dlaczego tak zdefiniowane ciało nie jest praktyczne}
Bardzo często (a nawet: w każdym sensownym przypadku) zdarza się, że w naszym wybranym ciele \(\mathbb{F}\) wielomian \(X^q - X\) nie ma dokładnie \(q\) pierwiastków. Wtedy możemy wykorzystać dopełnienie algebraiczne, czyli \(\overline{\mathbb{F}}\), w którym nasz wielomian będzie już miał odpowiadnią liczbę pierwiastków. Jednak na takim dopełnieniu algebraicznym jest trudno praktycznie pracować - ciężko jest konstrukcywnie opisać \(\overline{\mathbb{F}}\) aby móc robić na jego elementach operacje.
